\paragraph{Surjectivity from the deep level}
Write \(\widetilde Y_0^0=Y_0\perp\langle1\rangle\) for \(Y_0\) with one
positive line added. The deep level has modulus \(\pi_M^b\).
We first find a real point \(\omega\) for which the restrictions
\(g\mapsto\phi((g\oplus1)\omega)\) detect every form at the principal level
by its inner products.

Compare the lowest \(K\)-type weights at an infinite place.
With the sign reversal in the parameters, the common positive block
has weight zero. On the other entries the signs alternate. The weights
on one side decrease strictly in \([-3\kappa_0,-\kappa_0]\) and those of
the other are positive in the opposite order, where
\(\kappa_0=\kappa(\tau)+\Delta/2\).
The correction from the Harish--Chandra parameter to the lowest
\(K\)-type counts the entries above and below each entry with signs.
It is bounded by the rank. On the common positive block this count
cancels the original entry. On the remaining entries the gaps exceed
the correction and keep their ordering.

On the alternating block, let \(\mu_i\) be the negative weights on the
larger side and \(\mu'_j\) the positive weights on the smaller side.
Let \(P,Q\) be the respective monic polynomials with these roots.
For a polynomial \(F\) set
\[
 [F]_i=\frac{F(\mu_i)}{|P'(\mu_i)Q(\mu_i)|^{1/2}} .
\]
The Hermitian pairing of such real vectors is \((-1)^{\deg Q}\) times
the sum of the residues of \(F_1F_2/(PQ)\) at the \(\mu_i\).
As \(\deg P=\deg Q+1\) and \(\deg Q\) is even, \([Q]\) has Hermitian
length \(1\), as required for the extra positive line.
Write \(Q_j=Q/(X-\mu'_j)\). The residue of \(Q_j/P\) at infinity is zero,
so \([Q_j]\) is orthogonal to \([Q]\).
For \([Q_j]\) paired with \([Q_{j'}]\), the only other possible residue
is at \(\mu'_j\) when \(j=j'\). The resulting pairing is zero for
\(j\ne j'\) and \((-1)^{\deg Q+1}Q'(\mu'_j)/P(\mu'_j)\) on the diagonal.
Normalize these columns to Hermitian lengths of modulus \(1\) and
use the identity on the common positive block to make an isometry \(i_0\).
Because \((X-\mu'_j)Q_j=Q\), compression of
\(\operatorname{diag}(\mu_i)\) to the orthogonal complement of the
extra line gives \(\operatorname{diag}(\mu'_j)\).
Thus with \(w=(i_0,[Q])\) and \(g'=w(g\oplus1)w^{-1}\) the difference
of the lowest \(K\)-type phases
\(\operatorname{tr}(\mu(g'-I))-\operatorname{tr}(\mu'(g-I))\)
is zero.
\leanlinks{\leansrc{New/Determinant/AuxiliaryForms/DeepLevelPoincarePeriod}{Counting the entries of the lowest K-type weights}
  \leanlinksep \leansrc{New/Appendix/DiscreteSeries/StationaryPosition}{Constructing the line and its compression by residues}}

The entries are separated by at least \(1\) and their absolute values
are at most \(3\kappa_0\). Hence the coordinates of the line and the
isometry are bounded by \((6\kappa_0)^{M_0+1}\), where \(M_0\) is the
length of the smaller alternating block. Therefore
\(\|w\|^2\le(r+1)^4(6\kappa_0)^{2M_0+2}\le\beta\), with
\(\beta=(\sqrt M)^{\lceil x\rceil}\).
Pair the discrete series coefficients in this position.
Put \(D_0=(2r+1)^2\) and let \(J_g\) be the direct sum of the two group
matrices. In the Flensted--Jensen integrals the minor powers give decay
away from the compact subgroups. Near the compact subgroups, expand
\(\log(1+z)\) to order two. The linear phases cancel by compression
and the quadratic Gaussian integral gives a positive main term.
For the phase-aligned product \(B(g)\) and
\(D(g)=(2/\|J_g\|)^{9D_0}\) this gives
\[
 |B(g)|\le D(g),\qquad
 |B(g)|\le 2\RePart B(g)+6\beta^{-50D_0}D(g).
\]
The margin in the gaps here is at least \(10^{30}(D_0\log\beta)^3\).
On \(\|J_g-I\|\le(200\beta^5)^{-1}\) the modulus of \(B(g)\) is at
least \(1/4\). The ball \(\|J_g\|\le\beta+1\) has volume at most
\(\beta^{16D_0}\) times that of this small ball. The integral of
\(D(g)\) is at most \(16^{9D_0}\) times the volume of the unit ball.
These estimates show that the integrated error does not cancel
the contribution of the small ball, and give
\(v\le\beta^{20D_0}|\int B(g)\,dg|\), with \(v\) the volume of the
unit ball of the smaller group. Use \(\omega_\tau=w^{-1}\) at each
infinite place and take their product for \(\omega\).
\leanlinks{\leansrc{New/Appendix/DiscreteSeries/GaussianShape}{The Gaussian approximation and its error}
  \leanlinksep \leansrc{New/Appendix/DiscreteSeries/Periods}{Volume comparison and the integral lower bound}
  \leanlinksep \leansrc{New/Appendix/DiscreteSeries/PositionGroup}{Specialization to the codimension-one position}}

Suppose a form \(f\) on \(Y_0\) is orthogonal to every restriction
and is bounded by \(S\). We prove \(|f|\le S/2\) at every adelic point.
For this argument write \(L=\beta^{20r^4(f_E+1)}\) with \(f_E=|\Phi|\).
Where the first minimum is below
\((L^{-2}/(M C_{\mathrm{cusp}}^2))^r\), cusp vanishing bounds \(f\)
by the difference with its translate by a real point within \(L^{-2}\)
of the identity. Here \(C_{\mathrm{cusp}}\) is the lattice constant
depending on the rank and different. The real reproducing kernel
and its small-translation estimate bound the difference by
\(2S(\beta/L)\int K_\mathrm{maj}\le S/8\), where
\(\int K_\mathrm{maj}\le L/(16\beta)\).

At a point \(g\) with first minimum at least the threshold, form
the Poincaré series centred at \(g\oplus1\), using the discrete
series coefficient on \(\widetilde Y_0^0\) at infinity and the
indicator of the deep level at the finite places. Its period against
\(f\) is zero by orthogonality.
Unfolding, the identity coset contributes a real convolution of \(f\)
times the level volume. The other cosets are distinguished by the
last column of the rational matrix. The level congruence and the
lower bound for the first minimum separate these columns at the
infinite places by \(L^2\). Since the base column has norm at most \(L\),
count the remaining columns in balls from radius \(L^2/2\) and
apply the inverse-power decay. This gives
\((10/L)^{r^4}\le10/L\).
The real convolution is thus bounded by \(10S/L\) times the
factor controlling the translated matrix coefficient.
Convolve again with the reproducing kernel of \(Y_0\).
By Fubini and discrete series reproduction the result is \(f(g)\)
times the product of the real integrals bounded below above.
Dividing by that product cancels the factors \(v\) and the
powers of \(\beta\) from the majorant estimate, giving
\(|f(g)|\le (10/(16\beta))S\le S/2\).
Halving the uniform bound repeatedly now gives \(f=0\).
\leanlinks{\leansrc{New/Determinant/AuxiliaryForms/DeepLevelPoincarePeriod}{Last-column separation and the identity term}
  \leanlinksep \leansrc{New/Determinant/AuxiliaryForms/DeepLevelVanishing}{The cusp estimate and reproduction by real integrals}}

The forward chain of the second triangle at the deep level is
injective by the same induction on length using the exchanges of
Section~\ref{sec:exchange-proof}. Its transposed chain is surjective.
Indeed an \(f\) orthogonal to the image can be lifted by the forward
chain to \(Uf\), and transposition gives \(\|Uf\|^2=0\).
Injectivity gives \(f=0\).
Compose the transposed chain with restriction at \(\omega\) and
orthogonal projection onto the principal-level space. Since the
orthogonal complement to the image is zero, this gives the desired
surjection \(s\) onto the principal-level space.
\leanlinks{\leansrc{New/Determinant/AuxiliaryForms/FirstPath}{Surjectivity of the transposed chain and the restriction map}}
